\addchap{Notation}

One notation is used throughout the book. Where a chapter needs a local symbol that clashes with this list, the chapter says so; where a source paper uses a different notation, the chapter's notes translate it.

\newcommand{\notationrow}[2]{#1 & #2\\}
\newcommand{\notationhead}[1]{\multicolumn{2}{@{}l}{\textsc{\color{sea}#1}}\\[0.3ex]}
\noindent{\small
\renewcommand{\arraystretch}{1.1}%
\begin{tabular}{@{}>{\raggedright\arraybackslash}p{30mm}>{\raggedright\arraybackslash}p{86mm}@{}}
\notationhead{sizes}
\notationrow{$n$}{number of training samples}
\notationrow{$d$}{input dimension}
\notationrow{$m$}{width of a hidden layer (number of hidden units)}
\notationrow{$p$}{number of features or parameters of a linear model (the prime modulus in \cref{ch:grokking})}
\notationrow{$K$}{number of classes}
\notationrow{$N$}{context length (\cref{ch:icl,ch:induction}); model size, with $D$ the data size (\cref{ch:scaling})}
\notationrow{$L$, $\ell$}{depth, and the index of a layer}
\notationrow{$k$}{a degree, a sparsity, or an information exponent}
\addlinespace
\notationhead{data and models}
\notationrow{$x$, $y$}{input and label}
\notationrow{$f(x;\theta)$}{the predictor, with parameters $\theta$}
\notationrow{$w$, $W$}{a weight vector and a weight matrix}
\notationrow{$\phi$}{an activation function}
\notationrow{$w_\star$, $f_\star$, $W_\star$}{a subscript star marks the target or the teacher}
\notationrow{$\lambda^\star$, $\Gamma^\star$}{a superscript star marks an optimal value}
\notationrow{$r=\ip{w,w_\star}$}{the overlap of two unit vectors; $Q$, $R$, $T$ are matrices of overlaps}
\notationrow{$\sigma$}{the standard deviation of label noise}
\notationrow{$\sigma_w^2$, $\sigma_b^2$}{the variances of weights and biases at initialization}
\notationrow{$\alpha$}{the scale of initialization}
\addlinespace
\notationhead{training}
\notationrow{$\Loss(\theta)$, $\Lhat(\theta)$}{population loss and training loss}
\notationrow{$\eta$}{step size}
\notationrow{$t$}{time: continuous for flows, an integer for iterations}
\notationrow{$\lambda$}{ridge parameter or weight decay; eigenvalues always carry a subscript, as in $\lambda_k$, $\lmax$, $\lmin$}
\end{tabular}

\medskip\noindent
\begin{tabular}{@{}>{\raggedright\arraybackslash}p{30mm}>{\raggedright\arraybackslash}p{86mm}@{}}
\notationhead{probability and functions}
\notationrow{$\cN(\mu,\Sigma)$, $\Unif$}{Gaussian and uniform laws}
\notationrow{$\Sph^{d-1}$}{the unit sphere in $\R^d$}
\notationrow{$\He_k$, $h_k$}{probabilists' Hermite polynomials and their normalization $h_k=\He_k/\sqrt{k!}$, orthonormal under $\cN(0,1)$}
\notationrow{$\hat g_k$}{the Hermite coefficient $\E[g(z)h_k(z)]$, $z\sim\cN(0,1)$}
\notationrow{$\ind\{\cdot\}$}{indicator}
\addlinespace
\notationhead{linear algebra}
\notationrow{$\ip{u,v}$, $\norm{u}$}{Euclidean inner product and norm}
\notationrow{$\norm{A}_{\op}$, $\norm{A}_F$}{operator and Frobenius norms}
\notationrow{$A^\top$, $I_d$, $\mathbf 1_d$}{transpose, identity, all-ones vector}
\addlinespace
\notationhead{asymptotics}
\notationrow{$a\lesssim b$, $a\asymp b$}{$a\le Cb$, and $a\lesssim b\lesssim a$, for a constant $C$ depending only on parameters declared fixed}
\notationrow{$O$, $\Omega$, $\Theta$, $o$, $\tilde O$}{the usual order symbols; $\tilde O$ hides polylogarithmic factors}
\addlinespace
\notationhead{status of an argument}
\notationrow{$\blacksquare$}{end of a complete proof}
\notationrow{proof sketch}{an argument that omits steps, with a pointer to where they are}
\notationrow{$\Diamond$}{end of a heuristic derivation: an argument that is not a proof}
\notationrow{observation}{an empirical regularity, with its evidence}
\end{tabular}\par}
